\chapter{}
\label{chap:2}

\section{Умова}

Довести, що матриці Адамара задовольняють рекурентне співвідношення:
\begin{equation*}
    H_{n} = \begin{pmatrix}
        H_{n-1} &  H_{n-1} \\
        H_{n-1} & -H_{n-1} \\
    \end{pmatrix}, \quad n = 2, 3, \dots, \text{ де } 
        H_{1} = \begin{pmatrix} 1 & 1 \\ 1 & -1 \end{pmatrix}
\end{equation*}

\section{Розв'язання}
Дане рекурентне співвідношення (конструкція Сільвестра) є власне означенням матриць Уолша-Адамара для розмірності 
$N = 2^n$.
\begin{proof}
    ~\par Щоб довести, що побудована таким чином матриця $H_n$ є матрицею Адамара, треба перевірити головну властивість 
    матриця має складатися лише з $\pm 1$, а її рядки мають бути попарно ортогональними. Формулою це можна записати 
    так: $H_{n} \cdot H_{n}^{T} = 2^n \cdot I_{n}$, де $I$ це одинична матриця розмірності $n$. Найпростіше довести 
    за ММІ (не факультетом, а методом математичної індукції):
    \begin{enumerate}
        \item База індукції ($n = 1$):
            \begin{equation*} H_1 \cdot H_1^T = \begin{pmatrix} 1 & 1 \\ 1 & -1 \end{pmatrix} \cdot 
                \begin{pmatrix} 1 & 1 \\ 1 & -1 \end{pmatrix} = \begin{pmatrix} 2 & 0 \\ 0 & 2 \end{pmatrix} = 2 \cdot I_{2}
            \end{equation*}
            Властивість виконується.
        \item Припустимо, що для матриці $H_{n-1}$ (розмірності $2^{n-1} \times 2^{n-1}$) властивість теж виконується:
            \begin{equation*}
                H_{n-1} \cdot H_{n-1}^{T} = 2^{n-1} I_{2^{n-1}}
            \end{equation*}
            Варто зазначити, що матриці Адамара такого вигляду побудовою є симетричними, тобто $H_{n-1}^{T} = H_{n-1}$.
        \item Крок індукції ($n$): Перевіримо властивість для $H_{n}$, використовуючи блочне множення матриць:
            \begin{equation*}
                H_n \cdot H_n^T = \begin{pmatrix} H_{n-1} & H_{n-1} \\ H_{n-1} & -H_{n-1} \end{pmatrix} \cdot 
                \begin{pmatrix} H_{n-1}^T & H_{n-1}^T \\ H_{n-1}^T & -H_{n-1}^T \end{pmatrix}
            \end{equation*}
            Маємо:
            \begin{itemize}
                \item Лівий верхній блок: $H_{n-1}H_{n-1}^T + H_{n-1}H_{n-1}^T = 2(2^{n-1}I) = 2^n I$
                \item Правий верхній блок: $H_{n-1}H_{n-1}^T - H_{n-1}H_{n-1}^T = 0$
                \item Лівий нижній блок: $H_{n-1}H_{n-1}^T - H_{n-1}H_{n-1}^T = 0$
                \item Правий нижній блок: $H_{n-1}H_{n-1}^T + (-H_{n-1})(-H_{n-1}^T) = 2(2^{n-1}I) = 2^n I$
            \end{itemize}
            Зібравши докупи маємо:
            \begin{equation*}
                H_n H_n^T = \begin{pmatrix} 2^n I & 0 \\ 0 & 2^n I \end{pmatrix} = 2^n I_{2^n}
            \end{equation*}
        \end{enumerate}
        Отже, матриці задані цим рекурентним співвідношенням, є матрицями Адамара.
\end{proof}